%% ---------------------------------------------------------------------------
\section{Carrying the signer}
\label{sec:carrying}

If certificates carry the cost, the question is how many of them have to
travel. Table~\ref{tab:carriage} measures four answers under two placements
of the trust anchor.

\begin{table*}[t]
\caption{Carrying the signer. Manifest bytes by carriage form and algorithm, three-certificate hierarchy. \emph{Path} is the number of signatures path validation checks, set by where the trust anchor sits; \emph{Holds} is what the verifier must already have; \emph{int.} is the intermediate. $^a$C2PA~2.5 says the root should not be included. $^b$A reading the specification does not exclude, accepted by the reference implementation when the leaf's issuer is a configured anchor; clarification requested~\cite{c2paspec142}. Otherwise the leaf-only form does not conform.}
\label{tab:carriage}
\footnotesize
\setlength{\tabcolsep}{3.5pt}
\begin{tabular}{llclrrrrrr}
\toprule
 & & & &  &  &  & \multicolumn{3}{c}{ML-DSA}\\\cmidrule(l){8-10}
\texttt{x5chain} & Holds & Path & Conf. & Ed25519 & ES256 & PS256 & 44 & 65 & 87\\
\midrule
Leaf, int., root & root & 2 & yes$^a$ & 2,981 & 3,155 & 4,535 & 16,339 & 21,815 & 29,007\\
Leaf, int. & root & 2 & yes & 2,544 & 2,658 & 3,642 & 12,238 & 16,185 & 21,419\\
Leaf (issuer is anchor) & issuing CA & 1 & yes$^b$ & 1,849 & 1,906 & 2,495 & 7,875 & 10,293 & 13,569\\
Leaf (int.\ provisioned) & root, int. & 2 & no & 1,849 & 1,906 & 2,495 & 7,875 & 10,293 & 13,569\\
Thumbprint & root, int., leaf & 2 & no & 1,185 & 1,185 & 1,379 & 3,543 & 4,432 & 5,750\\
\midrule
\multicolumn{4}{l}{Saving, leaf only (\%)} & 27 & 28 & 31 & 36 & 36 & 37\\
\multicolumn{4}{l}{Saving, thumbprint (\%)} & 53 & 55 & 62 & 71 & 73 & 73\\
\multicolumn{4}{l}{Cost of the root (\%)} & 17 & 19 & 25 & 34 & 35 & 35\\
\bottomrule
\end{tabular}
\end{table*}


\paragraph{Leaf, intermediate and root} What the Adobe chains in the public
test corpus carry. The specification says the anchor's certificate should not
be included, and it costs \NcostRootDSAsixtyfive\% more than leaving it out
under ML-DSA-65 while changing no trust decision, since the verifier must
hold the root anyway. Dropping it is free.

\paragraph{Leaf and intermediate} The form C2PA~2.5 requires and the baseline
of every comparison in this paper.

\paragraph{Leaf only} The specification requires every intermediate CA certificate and says the trust anchor's should be left out. When the CA that issued the
leaf is itself on the trust list, it is the anchor, there is no intermediate
to carry, and the leaf alone is the conformant form. The C2PA trust list is
built to allow this: \NtlIssuing{} of its \NtlAnchors{} anchors are issuing
CAs rather than roots, among them Google's capture and media CAs, DigiCert's
claim-signing CAs and Adobe's product CA~\cite{c2patrustlist}. In that case
the manifest is \NsaveLeafDSAsixtyfive\% smaller under ML-DSA-65, and path
validation checks one signature instead of two, because an anchor's own
signature is never checked. When the leaf's issuer is not an anchor, carrying
the leaf alone works only for a verifier that already holds the intermediate,
which no specification provides for; Table~\ref{tab:validation} shows the
reference implementation returning \emph{Valid} rather than \emph{Trusted} in
exactly that case.

\paragraph{Thumbprint} The \texttt{x5t} header and no certificate, a further
\NsaveThumbDSAsixtyfive\% saving against leaf and intermediate. It is not
conformant, and the reference implementation reads it as a missing
credential.

Section~\ref{sec:prodchains} measures how often the conformant leaf-only
condition holds for production signers today: \NcorpProdListed{} of the
\NcorpProdChains{} distinct production chains we found end at a trust-list
anchor, and for \NcorpProdIssuer{} of them the leaf's issuer is the anchor. Our recommendation follows from both
measurements. Generators should never carry the root. A CA that issues
post-quantum signing certificates should have its issuing CA listed as a
trust anchor, as several already do for classical ones, and generators under
such a CA should carry the leaf alone; the saving is a third of the manifest
and one of two path signatures, and it needs no change to any specification.

\subsection{What each option gives up}
\label{sec:tradeoffs}

\emph{The trust list does not resolve a thumbprint.} It holds anchors, not
signers' leaf certificates, so a verifier holding only an \texttt{x5t} has
nowhere specification-defined to fetch the leaf from, and some signers issue
a leaf per image~\cite{google2025pixel}, for which there could be no such
directory.

\emph{Offline verification of an unknown signer.} A verifier with no network
that has never seen the signer can verify leaf and intermediate, can verify
the leaf alone when the leaf's issuer is an anchor it holds, and cannot verify
a thumbprint. Archival readers opening a file decades later are the case where
provenance matters most and fetching is least possible.

\emph{Privacy.} Remote thumbprint resolution reports every verification to
whoever hosts the resolver: which verifier looked up which signer, and when.
For a journalist checking a document, or a reader of politically sensitive
media, that is a disclosure the embedded forms do not make.

\subsection{Where verification time goes}
\label{sec:verifytime}

\begin{table*}[t]
\caption{Where the cryptographic work of validating one timestamped image goes, in milliseconds on Apple M1 Pro. \emph{Path} checks leaf and intermediate up to the root anchor; \emph{TST} and \emph{TSA path} are the timestamp's signature and certificate path; \emph{Hash} is SHA-256 over the content, for the 13th of the 24 Kodak images by size (113,540\,B) and for a 4\,MiB phone JPEG. Assertion hashing, ingredients, revocation and policy checks are outside the boundary. \emph{Total} is the sum of the row for each image. Times are medians over 10 independent runs, each the median of its own repetitions; Table~\ref{tab:ci} in the supplement gives their intervals across runs.}
\label{tab:verify}
\small
\setlength{\tabcolsep}{4pt}
\begin{tabular}{lrrrrrrrrr}
\toprule
 & Decode & Path & Manifest & TST & TSA & \multicolumn{2}{c}{Hash} & \multicolumn{2}{c}{Total}\\
\cmidrule(lr){7-8}\cmidrule(lr){9-10}
Algorithm & & & signature & & path & median & 4\,MiB & median & 4\,MiB\\
\midrule
Ed25519 & 0.001 & 0.181 & 0.081 & 0.080 & 0.180 & 0.046 & 1.70 & 0.571 & 2.23\\
ES256 & 0.001 & 0.148 & 0.056 & 0.054 & 0.148 & 0.046 & 1.71 & 0.456 & 2.12\\
PS256 & 0.001 & 0.069 & 0.021 & 0.016 & 0.069 & 0.046 & 1.71 & 0.223 & 1.89\\
ML-DSA-44 & 0.002 & 0.170 & 0.043 & 0.070 & 0.171 & 0.046 & 1.71 & 0.501 & 2.16\\
ML-DSA-65 & 0.002 & 0.251 & 0.068 & 0.108 & 0.251 & 0.046 & 1.70 & 0.729 & 2.39\\
ML-DSA-87 & 0.003 & 0.380 & 0.110 & 0.171 & 0.380 & 0.046 & 1.71 & 1.089 & 2.76\\
\bottomrule
\end{tabular}
\end{table*}


Certificates dominate verification for the same reason they dominate bytes,
but how much depends on what is counted. For the manifest signature and the
signer's path alone, the measurement most comparable to a TLS handshake, the
path is \NvPathShareSigDSAsixtyfive\% of the ML-DSA-65 work in
Table~\ref{tab:verify}, against the two thirds that two path signatures and
one manifest signature would suggest; the excess is that the path goes
through OpenSSL's X.509 stack, whose ML-DSA verification is slower than
liboqs's, and a verifier cannot choose a faster one for certificates because
liboqs has no ASN.1 layer. That fraction is conditional on the hierarchy: with
the issuing CA as anchor the path is one signature, and with a cached signer
(Section~\ref{sec:cache}) it is none.

Complete validation of a timestamped image adds the timestamp's signature and
path and the content hash, and Figure~\ref{fig:latency} shows the result. For a mid-sized Kodak image, the 13th of 24 by size, the signer's path falls to \NvPathShareDSAsixtyfive\% of
the cryptographic work under ML-DSA-65; for a 4\,MiB phone JPEG, hashing the
content dominates and the algorithm barely shows, \NvCompletePhonePS\,ms
under PS256 against \NvCompletePhoneDSAsixtyfive\,ms under ML-DSA-65. The
reference implementation, which also parses JUMBF and checks every assertion,
takes \NvalVerifyLo{} to \NvalVerifyHi\,ms per timestamped image
(Section~\ref{sec:patch}).

\begin{figure}[t]
\centering
\includegraphics[width=\linewidth]{figures/latency}
\Description{Horizontal stacked bars, one per algorithm, showing the time to
validate one timestamped image split into decode, signer's path, manifest
signature, timestamp signature, TSA path and content hash for a mid-sized Kodak
image. Totals run from about 0.2 to 1.1 milliseconds, and the signer's path and the TSA path are the largest segments.}
\caption{The cryptographic work of validating one timestamped mid-sized Kodak
image (the 13th of 24 by size) on the M1 Pro (Table~\ref{tab:verify}). Assertion hashing, ingredients,
revocation and policy checks are outside the boundary.}
\label{fig:latency}
\end{figure}

\subsection{Caching a validated signer}
\label{sec:cache}

A verifier that sees many manifests can cache certificates it has already
validated and skip path validation for them. Whether that pays depends on how
signers use certificates, so Table~\ref{tab:cache} in the supplement runs a
verifier over a stream from \Ncachepublishers{} publishers, each with its own
issuing CA, under three leaf lifecycles. With one long-lived leaf per
publisher, a cache of 4,096 certificates makes ML-DSA-65 verification
\NcacheLongDSAsixtyfive{} times faster. With a new leaf every hundred manifests the gain is about the same, \NcacheRotDSAsixtyfive{} times. With a new leaf for
every image, as signers that issue a certificate per capture do, only the
issuing CA is ever found in the cache, and the gain falls to
\NcachePerDSAsixtyfive{} times.

A cache also has to forget. Half-way through each stream we removed one
publisher's issuing CA from trust. A cache whose entries carry the trust-list
snapshot they were validated under, as C2PA~2.5 now makes possible, rejected
every later manifest from that publisher. A cache without that tag went on
accepting \NnaiveLongDSAsixtyfive{} of the \NnaiveTotLongDSAsixtyfive{} that
followed,
because their certificates had been validated before the change. A snapshot tag covers only changes to the trust list. A cached validation also goes stale when a certificate on the path expires or is revoked while the list stays the same, so we ran the same stream over thirty simulated days with the most frequent publisher's signing certificate expiring, or revoked, half-way (Table~\ref{tab:cachestale} in the supplement). A cache keyed on the certificate alone accepted \NstaleExpNaive{} of the \NstaleExpTot{} manifests signed after expiry and \NstaleRevNaive{} of \NstaleRevTot{} after revocation; checking each entry against its path's notAfter stops the first and not the second. What works is to key an entry on the certificate, the trust-list snapshot and the validation policy, and to let it live only while every certificate on its path is valid and the revocation status it was validated with is current. With seven-day OCSP responses such a cache still accepted \NstaleRevBoundSeven{} manifests, the last \NstaleRevBoundSevenH{} hours after revocation, and with one-day responses \NstaleRevBoundOne{}, the last after \NstaleRevBoundOneH{} hours, which is the window any verifier that trusts a response until its nextUpdate has. A verifier that validates at a historical time, from a trusted timestamp, needs that time in the key as well.

\subsection{Verification across cores}
\label{sec:scale}

A platform checking credentials on upload sees manifests from many publishers
and asks how many cores that takes. We verify a stream from 500 publishers,
drawn from a Zipf distribution, so that a handful of busy publishers send a large share and most send a few each, on every core of the M1 Pro,
the Xeon VM and the Pi~5, with and without a cached signer
(Table~\ref{tab:scale} and Figure~\ref{fig:scale}).

\begin{table}[t]
\caption{Manifests verified per second, in thousands, on one core and on every core of each machine. A stream from 500 publishers, drawn Zipf(1), each with a three-certificate chain. \emph{Uncached} validates the embedded chain for every manifest; \emph{cached} validates each publisher's chain once per worker and then checks only the signature, through OpenSSL or, for ML-DSA, liboqs. x5chain carries leaf and intermediate; the root comes from the verifier's store. Medians of three independent runs of 20,000 verifications each.}
\label{tab:scale}
\begin{tabular}{lrrrrrr}
\toprule
 & \multicolumn{2}{c}{Uncached} & \multicolumn{2}{c}{Cached} & \multicolumn{2}{c}{Cached, liboqs}\\
\cmidrule(lr){2-3}\cmidrule(lr){4-5}\cmidrule(lr){6-7}
Algorithm & 1 core & all & 1 core & all & 1 core & all\\
\midrule
\multicolumn{7}{l}{\emph{Apple M1 Pro, all = 8 cores}}\\
Ed25519 & 3.6 & 21.1 & 11.2 & 54.6 & -- & --\\
ES256 & 4.4 & 25.7 & 15.4 & 73.1 & -- & --\\
PS256 & 9.5 & 55.6 & 40.4 & 171.3 & -- & --\\
ML-DSA-44 & 3.8 & 21.9 & 11.2 & 53.3 & 19.6 & 87.4\\
ML-DSA-65 & 2.5 & 14.7 & 7.4 & 36.1 & 12.6 & 56.9\\
ML-DSA-87 & 1.6 & 9.6 & 4.8 & 23.1 & 7.9 & 35.6\\
\midrule
\multicolumn{7}{l}{\emph{Intel Xeon VM, all = 2 cores}}\\
Ed25519 & 2.4 & 4.6 & 7.4 & 13.9 & -- & --\\
ES256 & 2.7 & 5.2 & 10.1 & 18.5 & -- & --\\
PS256 & 6.4 & 12.9 & 26.6 & 47.8 & -- & --\\
ML-DSA-44 & 2.1 & 4.2 & 6.5 & 12.4 & 20.0 & 32.6\\
ML-DSA-65 & 1.5 & 2.9 & 4.2 & 7.9 & 13.2 & 24.1\\
ML-DSA-87 & 1.0 & 1.9 & 2.8 & 5.3 & 9.8 & 16.6\\
\midrule
\multicolumn{7}{l}{\emph{Raspberry Pi 5, all = 4 cores}}\\
Ed25519 & 2.1 & 8.3 & 6.5 & 23.1 & -- & --\\
ES256 & 1.3 & 5.3 & 4.6 & 15.9 & -- & --\\
PS256 & 2.2 & 8.9 & 7.9 & 27.0 & -- & --\\
ML-DSA-44 & 1.3 & 5.1 & 3.9 & 13.7 & 8.5 & 25.8\\
ML-DSA-65 & 0.9 & 3.5 & 2.6 & 9.3 & 5.6 & 17.5\\
ML-DSA-87 & 0.6 & 2.3 & 1.7 & 6.1 & 3.6 & 11.2\\
\bottomrule
\end{tabular}
\end{table}


\begin{figure}[t]
\centering
\includegraphics[width=\linewidth]{figures/scale}
\Description{Two line charts of ML-DSA-65 verifications per second against
worker processes, uncached and cached, with one line each for the M1 Pro, the
Xeon VM and the Raspberry Pi 5. Throughput rises close to linearly on the
Pi's four cores and the M1 Pro's six performance cores.}
\caption{ML-DSA-65 verification throughput against worker processes on three
platforms, without and with a cached signer (Table~\ref{tab:scale}).}
\label{fig:scale}
\end{figure}

Without a cache, one core verifies \NscMUncDSAsixtyfive{} ML-DSA-65 manifests
a second on the M1 Pro, \NscXUncDSAsixtyfive{} on the Xeon VM and
\NscPUncDSAsixtyfive{} on the Pi~5. On the M1 Pro that is \NscMVsES\% of the
ES256 rate and \NscMVsPS\% of the PS256 rate; the Xeon (\NscXVsES\% and
\NscXVsPS\%) and the Pi (\NscPVsES\% and \NscPVsPS\%) give the same order.
The work parallelises:
\NscPMax{} workers on the Pi give \NscPSpeedUncDSAsixtyfive{} times one
worker's rate, and \NscMMax{} on the M1 Pro give \NscMSpeedUncDSAsixtyfive{}
times, its two efficiency cores adding less than a performance core. At one
core's rate, ten thousand ML-DSA-65 manifests a second need
\NscMCoresUncDSAsixtyfive{} M1 Pro cores, \NscXCoresUncDSAsixtyfive{} Xeon
vCPUs or \NscPCoresUncDSAsixtyfive{} Pi cores, against
\NscMCoresUncES{}, \NscXCoresUncES{} and \NscPCoresUncES{} for ES256. A
cached signer leaves only the manifest signature, \NscMCacDSAsixtyfive{} a
second per M1 Pro core through OpenSSL and \NscMOqsDSAsixtyfive{} through
liboqs, and cuts those counts to \NscMCoresOqsDSAsixtyfive,
\NscXCoresOqsDSAsixtyfive{} and \NscPCoresOqsDSAsixtyfive. These are throughputs for the manifest signature and the signer's path, the boundary of Table~\ref{tab:scale}. A complete validation of a timestamped median image, with the timestamp's signature and path and the content hash, is \NvCompleteFactorDSAsixtyfive{} times that work under ML-DSA-65 (Table~\ref{tab:verify}), and parsing, assertion checks and revocation lookups come on top; the core counts scale accordingly and are a lower bound for a full ingest pipeline. Within that boundary, post-quantum verification at upload scale is a matter of a few more cores, not of new hardware.
